% =============================================================================
%  CHAPTER 4 — IMPLEMENTATION
%  RUBRIC: Execution & Technical Quality (50%). BUDGET: ~8-12 pages.
%  DETAILED OUTLINE (facts from your README / PROJECT_CONTEXT.md / EVALUATION.md).
%  Expand into prose IN YOUR OWN WORDS. Go DEEP on the hard parts (the two
%  engineering challenges) — depth over breadth is what reaches the 70-84 band.
%  Include short, well-explained code listings (lstlisting), not whole files.
% =============================================================================
\chapter{Implementation}
\label{ch:impl}

\section{Plugin Structure and the Tool Server}
\label{sec:impl-server}
% OUTLINE:
%   - AgentToolServer.kt: starts an HTTP server on 127.0.0.1:6473 when a project opens.
%   - Endpoints: POST /tools (ToolsHandler dispatch), GET /tools/schema (SchemaHandler,
%     Claude-compatible tool defs), GET /status (StatusHandler).
%   - Dispatch: ToolsHandler maps {tool, params} -> a project service method.
%   - Threading: PSI mutations must run on the EDT inside a write action; describe
%     the runProcessorOnEdt dispatcher used by the refactoring processors.
%   - Key source files (one line each): RefactorService.kt, PsiCreationService.kt,
%     KotlinCreationService.kt, ProjectContextProvider.kt, ApiKeyService.kt.
\texttt{AgentToolServer} starts a \texttt{com.sun.net.httpserver.HttpServer} bound to
\texttt{127.0.0.1:6473} when a project opens (via a \texttt{StartupActivity}) and
registers three contexts: \texttt{POST /tools} (\texttt{ToolsHandler}),
\texttt{GET /tools/schema} (\texttt{SchemaHandler}, which serves the function-calling
tool definitions), and \texttt{GET /status} (\texttt{StatusHandler}). A request carries
a JSON body \texttt{\{tool, params\}}; \texttt{ToolsHandler.dispatch} maps the tool name
to a method on a project-scoped service --- \texttt{RefactorService},
\texttt{PsiCreationService}, or \texttt{KotlinCreationService}.

The one threading subtlety is that IntelliJ PSI mutations must run on the
event-dispatch thread (EDT) inside a write action, whereas the HTTP request arrives on
a server thread. Each mutating call is therefore marshalled through a small dispatcher,
\texttt{runProcessorOnEdt}, which first waits for indexing to finish
(\texttt{DumbService.waitForSmartMode}), then \texttt{invokeAndWait}s on the EDT,
committing open documents before running the operation and committing and saving all
documents afterwards so the change is flushed to disk before the HTTP response returns.
Simpler, non-processor writes use a sibling dispatcher that wraps the body in a
\texttt{WriteCommandAction}. This is what lets a stateless HTTP call drive a stateful,
thread-confined IDE engine safely.

\section{Refactoring Operations via IntelliJ PSI}
\label{sec:impl-refactor}
% OUTLINE:
%   - Each operation delegates to an IntelliJ processor:
%       rename        -> RenameProcessor
%       move class    -> MoveClassesOrPackagesProcessor
%       change sig    -> ChangeSignatureProcessor
%       safe delete   -> SafeDeleteProcessor
%       inline        -> InlineMethodProcessor
%   - WORKED EXAMPLE (pick rename): trace renameByQualifiedName from the HTTP
%     request, through resolveQualifiedName -> PsiElement, to RenameProcessor.run(),
%     to the updated files. Show the ~5-line core as a listing.
Every refactoring operation is a thin delegation to the corresponding IntelliJ
processor: rename to \texttt{RenameProcessor}, move to
\texttt{MoveClassesOrPackagesProcessor}, change-signature to
\texttt{ChangeSignatureProcessor}, safe-delete to \texttt{SafeDeleteProcessor}, and
inline to \texttt{InlineMethodProcessor}. The service contributes symbol resolution and
headless execution; the processor contributes the correctness guarantee.

A rename is representative. The HTTP layer calls
\texttt{RefactorService.renameByQualifiedName}, whose core is short:
\begin{lstlisting}[basicstyle=\ttfamily\footnotesize]
fun renameByQualifiedName(qualifiedName: String, newName: String, ...): Result {
    refreshProjectFromDisk()                  // reconcile VFS with on-disk state
    val element = ReadAction.compute {
        resolveQualifiedName(qualifiedName) as? PsiNamedElement
    } ?: return Result.Err("could not resolve \"$qualifiedName\"")
    return rename(element, newName, ...)      // runs the processor on the EDT
}
\end{lstlisting}
The qualified name is resolved to a \texttt{PsiNamedElement} inside a read action;
\texttt{rename} then constructs a \texttt{HeadlessRenameProcessor} --- a
\texttt{RenameProcessor} subclass that runs without the interactive preview and
conflict dialogs that would otherwise block an unattended IDE --- and calls
\texttt{run()} inside \texttt{runProcessorOnEdt}. The processor rewrites the declaration
and every resolved reference across the project as a single atomic refactoring before
the files are saved.

\section{Headless Extract-Method}
\label{sec:impl-extract}
% OUTLINE:
%   - Prior extract-method tools (e.g. EM-Assist) are editor-bound: they need an
%     open editor with a caret selection.
%   - This uses IntelliJ's newImpl MethodExtractor.extractMethod(ExtractOptions),
%     driven from a PsiFile + text range with no UI.
%   - Note why this is non-trivial (building ExtractOptions without an editor).
Extract-method is the one operation with no clean headless entry point. The
user-facing refactoring is editor-bound --- tools such as EM-Assist require an open
editor with a caret selection --- because the public API threads an \texttt{Editor}
through option-gathering and the confirmation UI. This plugin instead drives IntelliJ's
\texttt{newImpl} extraction pipeline directly from a file and a text range, with no
editor: \texttt{ExtractSelector().suggestElementsToExtract(psiFile, range)} identifies
the PSI elements covered by the selection, \texttt{ExtractMethodPipeline.findAllOptionsToExtract}
computes the candidate \texttt{ExtractOptions}, and the first option --- with its
\texttt{methodName} replaced by the requested name --- is handed to
\texttt{MethodExtractor().extractMethod(options)} on the EDT. The non-trivial part is
assembling a valid \texttt{ExtractOptions} (the extracted parameters, return type, and
target) without the editor that normally supplies that information; the \texttt{newImpl}
structures expose exactly the entry points needed to do so, where the legacy
extract-method handler does not.

\section{Symbol Creation (Java and Kotlin)}
\label{sec:impl-create}
% OUTLINE:
%   - Java creation via PsiElementFactory (+ auto-format): add_field/method/inner_class/file.
%   - Kotlin creation via KtPsiFactory: add_kt_property/function, create_kotlin_file.
%   - One schema, two languages, no duplication in the HTTP surface; Kotlin tools
%     load only when the Kotlin plugin is present.
Creation tools add new code rather than transform existing code, and are built on the
PSI element factories. The Java tools (\texttt{add\_field}, \texttt{add\_method},
\texttt{add\_inner\_class}, \texttt{create\_java\_file}) build elements with
\texttt{PsiElementFactory} and reformat them into the target class or file; the Kotlin
tools (\texttt{add\_kt\_property}, \texttt{add\_kt\_function},
\texttt{create\_kotlin\_file}) do the same through \texttt{KtPsiFactory}. Both languages
are exposed through one JSON tool schema, so an agent works against a single uniform
surface regardless of language. The Kotlin tools are resolved lazily and advertised only
when the Kotlin plugin is present, so the system degrades gracefully to a Java-only
surface when it is not, rather than failing.

\section{Project Context Injection}
\label{sec:impl-context}
% OUTLINE:
%   - ProjectContextProvider builds a system prompt each chat turn containing:
%     project name + source roots, open file(s) with full symbol listings
%     (names, kinds, signatures, offsets), tool-usage + qualified-name guidance.
%   - Effect: the agent is productive on a project it has never seen, with no
%     manual paste of file paths or offsets.
For an agent to act on a project it has never seen, it must discover what is there
without the user pasting file paths or offsets. \texttt{ProjectContextProvider} assembles
this automatically: on each turn it builds a context block containing the project name
and source roots, a symbol listing for the relevant file(s) --- names, kinds,
signatures, and offsets --- and short guidance on the tool surface and on addressing
symbols by qualified name. Because the agent receives the resolved symbol inventory up
front, it can issue a correct \texttt{rename\_symbol} or \texttt{find\_usages} call by
name on its first turn instead of spending turns reading files to locate a target,
which is what makes it productive on an unfamiliar codebase out of the box.

\section{Engineering Challenges}
\label{sec:impl-challenges}
% These two subsections are your strongest depth signal — write them fully.

\subsection{VFS Staleness After External File Changes}
\label{sec:impl-vfs}
% OUTLINE (from EVALUATION.md "Plugin robustness fix"):
%   - Symptom: initial structured runs scored 6/8 due to cross-task contamination.
%   - Root cause: git resets between tasks updated the DISK, but IntelliJ's
%     in-memory Documents retained the previous task's state; saveAllDocuments()
%     then silently failed (IntelliJ detected an external-change conflict and
%     refused to overwrite).
%   - Fix: every mutating entry point (renameByQualifiedName, safeDeleteByQualifiedName,
%     moveClass, changeSignature, inlineMethod, extractMethod, extractVariable) and
%     every PSI-creation entry point (addField, addMethod, addInnerClass) now calls a
%     synchronous refresh(false, true) on the source roots before reading PSI.
%   - Why it matters: robustness to external changes (git, external editors, CI) is a
%     property a production plugin needs anyway.
% --- AI FIRST DRAFT from your EVALUATION.md notes. Edit freely; then I polish. ---
The benchmark harness exposed a subtle but consequential interaction between
IntelliJ's in-memory model of a file and that file's state on disk. To keep each
benchmark task independent, the runner commits the project's baseline state
before a task and performs a \texttt{git reset} back to that commit afterwards,
so the next task begins from an identical starting point. Early structured runs
nevertheless scored only 6 of 8 tasks, and the failures were not reproducible in
isolation: a task that failed inside a full run passed when executed on its own.

The cause was a divergence between two representations of the same file. A
\texttt{git reset} rewrites the file on disk, but IntelliJ continues to hold an
in-memory \texttt{Document} reflecting the \emph{previous} task's edits. When the
plugin then applied a refactoring and called \texttt{saveAllDocuments()},
IntelliJ detected that the on-disk file had changed beneath its stale
\texttt{Document}, treated this as an external-modification conflict, and
silently declined to overwrite the file. The refactoring had executed correctly
against the PSI, but its result was never persisted --- so the validation step,
which compiles the project from disk, saw stale contents and reported a failure.

The fix makes every entry point defensive against external changes. Each mutating
operation (\texttt{renameByQualifiedName}, \texttt{safeDeleteByQualifiedName},
\texttt{moveClass}, \texttt{changeSignature}, \texttt{inlineMethod},
\texttt{extractMethod}, \texttt{extractVariable}) and each creation operation
(\texttt{addField}, \texttt{addMethod}, \texttt{addInnerClass}) now performs a
synchronous \texttt{refresh(false, true)} of the project's source roots before it
reads the PSI, forcing IntelliJ's virtual file system to reconcile with disk so
that all subsequent work proceeds from the current contents. With this change the
cross-task contamination disappeared and the suite scored 8 of 8.

Although the bug surfaced through the benchmark's reset-between-tasks pattern, the
underlying requirement is general: a plugin that exposes refactoring operations to
an external agent must tolerate the file system changing beneath it --- through
Git operations, an external editor, or a CI step --- rather than assuming it is
the sole writer. Refreshing before each operation is therefore not a
benchmark-specific workaround but a correctness property a production plugin needs
regardless.

\subsection{Class Rename Must Trigger File Rename}
\label{sec:impl-filerename}
% OUTLINE (from EVALUATION.md "Class rename now guarantees file rename"):
%   - Symptom: renaming a public class updated references in memory but did NOT
%     rename Foo.java -> Bar.java.
%   - Root cause: the initial implementation used manual PSI text replacement
%     (handleElementRename + setName), which skips IntelliJ's file-rename side effect.
%   - Fix: route through RenameProcessor(project, element, newName,
%     searchInComments, searchTextOccurrences).run() (the same runProcessorOnEdt
%     dispatcher used by the other processors), making the file rename part of the
%     same atomic refactoring.
%   - Tie to benchmark task pc-rename-002 (CrashController -> PanicController),
%     validated by compile_and_file_exists (PanicController.java exists).
% --- AI FIRST DRAFT from your EVALUATION.md notes. Edit freely; then I polish. ---
A second defect concerned the gap between renaming a \emph{symbol} and renaming
the \emph{file} that declares it. The initial implementation of class rename
manipulated the PSI directly, calling \texttt{handleElementRename}/\texttt{setName}
to change the class identifier and update its references in memory. For most
symbols this is sufficient; but a public Java class must reside in a file of the
same name, so renaming the class must also rename \texttt{Foo.java} to
\texttt{Bar.java}. The manual PSI path updated the declaration and its usages but
silently skipped this file-rename side effect, leaving a \texttt{PanicController}
class inside a file still called \texttt{CrashController.java} --- which does not
compile.

The fix routes class rename through IntelliJ's \texttt{RenameProcessor}
(\texttt{RenameProcessor(project, element, newName, searchInComments,
searchTextOccurrences).run()}), invoked through the same
\texttt{runProcessorOnEdt} dispatcher already used for moves and signature
changes. Because \texttt{RenameProcessor} models the rename at the level the IDE
itself uses, the file rename happens as part of the same atomic refactoring
rather than as a separate step the plugin must remember to perform.

This case is captured directly by benchmark task \texttt{pc-rename-002} (renaming
\texttt{CrashController} to \texttt{PanicController}), whose validation was
deliberately changed from a RefactoringMiner check to
\texttt{compile\_and\_file\_exists}: it asserts that \texttt{PanicController.java}
exists at the expected path after the operation, a direct and observable signal
that the file rename actually occurred. It also illustrates the project's central
thesis in miniature: delegating to the IDE's own refactoring engine yields a
correct, whole-operation result --- declaration, references, and file together
--- whereas the manual approach reproduced exactly the kind of partial,
inconsistent edit that a text-level tool would produce.
